% Reviewed translation: PDF pages 288--291 / printed pages 274--277.
% The original scan is authoritative; mathematical tokens were checked against it.
\MathBookSourcePage{288}
继续 Definition~\ref{def:incompressible-vector-potential-pressure-reference-space}.
令 $\kappa\subset\mathbb R^3$ 为四面体, 面记为 $f$,
$\mathbf u\in H^1(\kappa)^3$. 定义两组矩
\[
N_f(\mathbf u)=
\left\{\int_f\mathbf u\cdot\mathbf n\,q\,ds;
\ q\in P_{l-1}(f)\right\},
\qquad
N_\kappa(\mathbf u)=
\left\{\int_\kappa\mathbf u\cdot\mathbf q\,dx;
\ \mathbf q\in P_{l-2}^3(\kappa)\right\}.
\]
对 $\mathbf u\in D_l$, 每个面上
$\mathbf u\cdot\mathbf n\in P_{l-1}$.
此外, $D_l$ 中的无散向量场属于 $P_{l-1}^3$.

令 $\widehat\kappa$ 为单位参考四面体.
这里使用逆变变换
\begin{equation}\label{eq:incompressible-pressure-contravariant-transform-3d}
\widehat{\mathbf u}=B_\kappa^{-1}(\mathbf u\circ F_\kappa).
\end{equation}
它与二维逆变变换相差常数因子
$J_F=|\det B_\kappa|$, 并完全保持散度:
\begin{equation}\label{eq:incompressible-pressure-divergence-transform-3d}
(\operatorname{div}\mathbf u)\circ F_\kappa
=\operatorname{div}\widehat{\mathbf u}.
\end{equation}

\begin{Proposition}\label{prop:incompressible-pressure-reference-unisolvence-3d}
空间 $D_l$ 在
\eqref{eq:incompressible-pressure-contravariant-transform-3d} 下不变,
且 $\mathbf u$ 在 $\kappa$ 上的矩全为零, 当且仅当
$\widehat{\mathbf u}$ 在 $\widehat\kappa$ 上的矩全为零.
此外, $D_l$ 中的向量场由 $N_f,N_\kappa$ 唯一确定;
其在给定面 $f$ 上的法向分量只依赖于该面上的 $N_f$.
\end{Proposition}

\begin{Definition}\label{def:incompressible-pressure-interpolant-3d}
对任意四面体 $\kappa$ 和
$\mathbf u\in H^1(\kappa)^3$,
$\omega_\kappa\mathbf u$ 是 $D_l$ 中与 $\mathbf u$
具有相同矩的唯一多项式, 即
\[
N_\kappa(\mathbf u-\omega_\kappa\mathbf u)=0,
\qquad
N_f(\mathbf u-\omega_\kappa\mathbf u)=0.
\]
\end{Definition}
Proposition~\ref{prop:incompressible-pressure-reference-unisolvence-3d}
蕴含相应的仿射交换关系. 此外,
\[
\operatorname{div}(\mathbf u-\omega_\kappa\mathbf u)=0
\quad\text{于 }\kappa\quad\forall\mathbf u\in P_l^3,
\]
且 $\operatorname{div}\mathbf u=0$ 蕴含
$\operatorname{div}(\omega_\kappa\mathbf u)=0$.

\MathBookSourcePage{289}
设 $\Omega$ 为有界多面体, $\mathcal T_h$ 是其剖分族. 定义
\begin{equation}\label{eq:incompressible-pressure-spaces-3d}
\left\{
\begin{aligned}
D_h&=\{\mathbf u_h\in H(\operatorname{div};\Omega);
\ \mathbf u_h|_\kappa\in D_l
\ \forall\kappa\in\mathcal T_h\},\\
D_{0h}&=D_h\cap H_0(\operatorname{div};\Omega),\\
Q_h&=\{p_h\in L_0^2(\Omega);
\ p_h|_\kappa\in P_{l-1}
\ \forall\kappa\in\mathcal T_h\}.
\end{aligned}
\right.
\end{equation}
并定义全局插值
$(\omega_h\mathbf u)|_\kappa=\omega_\kappa\mathbf u$.
对每个 $\mathbf u_h\in D_{0h}$,
$\operatorname{div}\mathbf u_h\in Q_h$;
对每个 $\boldsymbol\mu_h\in\Phi_h$,
$\operatorname{curl}\boldsymbol\mu_h\in D_{0h}$.
若 $\mathbf u\in H^1(\Omega)^3$, 则
$\omega_h\mathbf u\in D_h$;
若此外 $\mathbf u\cdot\mathbf n=0$, 则
$\omega_h\mathbf u\in D_{0h}$.

\begin{Proposition}\label{prop:incompressible-pressure-approximation-3d}
令 $\mathcal T_h$ 是 $\overline\Omega$ 的正则剖分族,
$D_h$ 由 \eqref{eq:incompressible-pressure-spaces-3d} 定义.
则
\begin{equation}\label{eq:incompressible-pressure-approximation-3d}
\left\{
\begin{aligned}
\|\mathbf u-\omega_h\mathbf u\|_{0,\Omega}
&\leq C_1h^l|\mathbf u|_{l,\Omega}
&&\forall\mathbf u\in H^l(\Omega)^3,\\
\|\operatorname{div}(\mathbf u-\omega_h\mathbf u)\|_{0,\Omega}
&\leq C_2h^l|\mathbf u|_{l+1,\Omega}
&&\forall\mathbf u\in H^{l+1}(\Omega)^3.
\end{aligned}
\right.
\end{equation}
\end{Proposition}

\begin{Lemma}\label{lem:incompressible-pressure-curl-characterization-3d}
令 $\Omega$ 是有界多面体, 边界连通分支为
$\Gamma_i$, $0\leq i\leq p$.
$\mathbf u_h\in D_h$ (相应地 $D_{0h}$) 满足
\[
\operatorname{div}\mathbf u_h=0\text{ 于 }\Omega,
\qquad
\int_{\Gamma_i}\mathbf u_h\cdot\mathbf n\,ds=0
\quad(0\leq i\leq p),
\]
当且仅当存在 $\boldsymbol\phi_h\in M_h$
(相应地 $\Phi_h$) 使
$\mathbf u_h=\operatorname{curl}\boldsymbol\phi_h$.
\end{Lemma}

\begin{Proof}
一个方向已知. 反之, 连续向量势定理给出
$\boldsymbol\phi\in H^1(\Omega)^3$ 使
$\mathbf u_h=\operatorname{curl}\boldsymbol\phi$.
又 $\mathbf u_h\in H^\alpha(\Omega)^3$, $0<\alpha<1/2$,
蕴含 $\boldsymbol\phi\in H^{1+\alpha}(\Omega)^3$,
故 $r_h\boldsymbol\phi$ 有定义.
下面证明
$\mathbf u_h=\operatorname{curl}r_h\boldsymbol\phi$.

\MathBookSourcePage{290}
即需证明
$\operatorname{curl}(\boldsymbol\phi-r_h\boldsymbol\phi)=0$.
一方面, 该旋度在每个单元上属于 $P_{l-1}^3$;
另一方面,
\[
\int_\kappa\operatorname{curl}(\boldsymbol\phi-r_\kappa\boldsymbol\phi)
\cdot\mathbf q\,dx=0
\quad\forall\mathbf q\in P_{l-2}^3(\kappa),
\]
且其法向分量在每个面上为零.
如 Lemma~\ref{lem:incompressible-nedelec-tetrahedron-unisolvence}
的论证, 得它在每个单元上为零;
又因差属于 $H(\operatorname{curl};\Omega)$,
故在整个 $\Omega$ 上为零.

若 $\mathbf u_h\cdot\mathbf n=0$ 于 $\Gamma$,
还需使 $\boldsymbol\phi_h\times\mathbf n=0$.
沿连续向量势定理的证明, 取包含 $\overline\Omega$ 的开球 $\mathcal O$,
构造 $q\in H^2(\mathcal O)$ 使
$\operatorname{grad}q\times\mathbf n
=\boldsymbol\phi\times\mathbf n$ 于 $\Gamma$.
不要求 $\Omega$ 额外正则, 因为 $\operatorname{grad}q$ 不必无散.
于是 $r_h(\boldsymbol\phi-\operatorname{grad}q)$
就是 $\Phi_h$ 中所需向量势.
\end{Proof}

\begin{Remark}\label{rem:incompressible-pressure-unique-discrete-potential-3d}
由 Corollary~\ref{cor:incompressible-nedelec-discrete-decomposition},
对每个 $D_{0h}$ 中的无散向量场 $\mathbf u_h$,
存在唯一 $\boldsymbol\phi_h\in\Phi_{h0}$ 使
$\mathbf u_h=\operatorname{curl}\boldsymbol\phi_h$.
在 Proposition~\ref{prop:incompressible-nedelec-uniform-discrete-poincare}
的假设下,
\[
\|\boldsymbol\phi_h\|_{H(\operatorname{curl};\Omega)}
\leq C\|\mathbf u_h\|_{0,\Omega}.
\]
\end{Remark}

由该 Lemma 与向量势--涡量格式, 可写成同时求
$\mathbf u_h,\boldsymbol\omega_h,p_h$ 的问题:
求 $(\mathbf u_h,\boldsymbol\omega_h)
\in D_{0h}\times M_h$ 和 $p_h\in Q_h$, 使
\begin{equation}\label{eq:incompressible-pressure-full-discrete-problem-3d}
\left\{
\begin{aligned}
\nu(\operatorname{curl}\boldsymbol\omega_h,\mathbf v_h)
-(p_h,\operatorname{div}\mathbf v_h)
&=(\mathbf f,\mathbf v_h)
&&\forall\mathbf v_h\in D_{0h},\\
(\mathbf u_h,\operatorname{curl}\boldsymbol\mu_h)
&=(\boldsymbol\omega_h,\boldsymbol\mu_h)
&&\forall\boldsymbol\mu_h\in M_h,\\
\operatorname{div}\mathbf u_h&=0
&&\text{于 }\Omega.
\end{aligned}
\right.
\end{equation}
压力的存在唯一性由下列 inf--sup 结果保证;
Lemma~\ref{lem:incompressible-pressure-curl-characterization-3d}
又保证该问题与前两个离散问题等价.

\MathBookSourcePage{291}
\begin{Lemma}\label{lem:incompressible-pressure-inf-sup-3d}
令 $\Omega$ 为有界多面体. 对每个 $p_h\in Q_h$,
存在 $\mathbf v_h\in D_{0h}$ 和
$\boldsymbol\theta_h\in M_h$, 使
\[
\operatorname{div}\mathbf v_h=p_h\quad\text{于 }\Omega,
\qquad
(\boldsymbol\theta_h,\boldsymbol\mu_h)
=(\mathbf v_h,\operatorname{curl}\boldsymbol\mu_h)
\quad\forall\boldsymbol\mu_h\in M_h.
\]
若 $\mathcal T_h$ 正则, 则
\begin{equation}\label{eq:incompressible-pressure-inf-sup-3d}
\|\boldsymbol\theta_h\|_{0,\Omega}
+\|\mathbf v_h\|_{H(\operatorname{div};\Omega)}
\leq C\|p_h\|_{0,\Omega},
\end{equation}
其中 $C$ 与 $h,p_h$ 无关.
\end{Lemma}

\begin{Theorem}\label{thm:incompressible-pressure-final-error-3d}
令 $\Omega$ 是 $\mathbb R^3$ 中的有界多面体.
问题 \eqref{eq:incompressible-pressure-full-discrete-problem-3d}
在 $D_{0h}\times M_h\times Q_h$ 中有唯一解
$(\mathbf u_h,\boldsymbol\omega_h,p_h)$,
其中 $(\mathbf u_h=\operatorname{curl}\boldsymbol\psi_h,
\boldsymbol\omega_h)$ 是约束离散向量势问题的解.

若此外 Theorem~\ref{thm:incompressible-vector-potential-final-error}
的假设成立, 且精确压力 $p\in H^l(\Omega)$, $l\geq1$, 则
\begin{equation}\label{eq:incompressible-pressure-final-error-3d}
\|p-p_h\|_{0,\Omega}
\leq C\{h^l|p|_{l,\Omega}
+h^l\|\boldsymbol\omega\|_{l+1,\Omega}
+h^{l-1}|\boldsymbol\psi|_{l+1,\Omega}\},
\end{equation}
其中 $C$ 与 $h,p,\boldsymbol\omega,\boldsymbol\psi$ 无关.
\end{Theorem}
